%=============================================================================
%  Beamer slides: Theoretical Model (Section 4) of the Job Market Paper
%  "Random Subspace Method for Forecast Combination"
%  Motivation drawn from RSM_paper_v2; theory follows JMP Section 4.
%  Compile with pdflatex.
%=============================================================================
\documentclass[11pt,aspectratio=169]{beamer}

\usetheme{Madrid}
\usecolortheme{seahorse}
\setbeamertemplate{navigation symbols}{}
\setbeamertemplate{caption}[numbered]
\setbeamertemplate{itemize items}[circle]

\usepackage{amsmath,amssymb,amsthm}
\usepackage{bm}
\usepackage{bbm}
\usepackage{booktabs}
\usepackage{xcolor}
\usepackage[round,authoryear]{natbib}

%---- accent color -----------------------------------------------------------
\definecolor{rsmblue}{RGB}{31,73,125}
\setbeamercolor{frametitle}{bg=rsmblue,fg=white}
\setbeamercolor{title}{fg=rsmblue}
\setbeamercolor{structure}{fg=rsmblue}

%---- notation macros (consistent with the paper) ----------------------------
\newcommand{\Sig}{\Sigma_{e,m}}
\newcommand{\Sigk}{\Sigma_{e,k,S}}
\newcommand{\wopt}{w^{\mathrm{opt}}}
\newcommand{\wave}{w^{\mathrm{ave}}}
\newcommand{\Qopt}{Q^{\mathrm{opt}}}
\newcommand{\Qave}{Q^{\mathrm{ave}}}
\newcommand{\Qsub}{Q_k^{\mathrm{sub}}}
\newcommand{\Qbag}{Q_k^{\mathrm{bag},\infty}}
\newcommand{\QbagG}{Q_k^{\mathrm{bag},G}}
\newcommand{\Lopt}{L^{\mathrm{opt}}}
\newcommand{\Lave}{L^{\mathrm{ave}}}
\newcommand{\Lsub}{L_k^{\mathrm{sub}}}
\newcommand{\Lbag}{L_k^{\mathrm{bag},\infty}}
\newcommand{\Rk}{R_k^{\mathrm{bag}}}
\newcommand{\E}{\mathbb{E}}
\newcommand{\Var}{\mathrm{Var}}
\newcommand{\Cov}{\mathrm{Cov}}
\newcommand{\im}{\bm{\iota}_m}
\newcommand{\ik}{\bm{\iota}_k}

%---- theorem-like blocks ----------------------------------------------------
\newtheorem{assumptionX}{Assumption}
\newtheorem{propositionX}{Proposition}
\newtheorem{definitionX}{Definition}

%---- title ------------------------------------------------------------------
\title[RSM: Theoretical Model]{The Random Subspace Method for Forecast Combination}
\subtitle{Theoretical Model}
\author[Kozyrev]{Boris Kozyrev}
\institute[]{Job Market Paper}
\date{\today}

\begin{document}

%=============================================================================
\begin{frame}[plain]
  \titlepage
\end{frame}

%=============================================================================
\begin{frame}{Roadmap}
  \tableofcontents
\end{frame}

%=============================================================================
\section{Motivation}
%=============================================================================

\begin{frame}{The forecast-combination puzzle}
  \begin{block}{A long-standing puzzle}
    The \emph{equal-weighted mean} of individual forecasts typically
    \textbf{outperforms} the estimated minimum-variance ``optimal''
    combination --- even though the latter solves a simple constrained
    OLS problem.
    \hfill {\footnotesize \citep*{BatesGranger1969,GrangerRamanathan1984}}
  \end{block}

  \textbf{Two standard explanations:}
  \begin{itemize}
    \item \textbf{Estimation error.} Optimal weights must be estimated;
      with $m$ forecasts and short samples $T$, the noise overwhelms the gains.
    \item \textbf{Small population gains.} \citet*{ElliottLiao2025}: even
      \emph{before} any finite-sample noise, the population gain
      $\Qave-\Qopt$ is often tiny under empirically plausible covariance
      structures (probably, not true in the endogenous case).
  \end{itemize}

  \begin{alertblock}{This paper}
    Both forces point to using \emph{fewer} forecasts. The Random Subspace
    Method (RSM) tries to answer how many of them are needed.
  \end{alertblock}
\end{frame}

%-----------------------------------------------------------------------------
\begin{frame}{The inflation factor}
  Estimating weights inflates the population loss. For the full
  $m$-weight optimal combination (constrained OLS):
  \[
    I_m(T)=\frac{T-1}{T-m}\approx 1+\frac{m-1}{T}.
  \]

  \citet*{ElliottLiao2025} call $(m-1)/T$ the \emph{best case} estimation
  cost. Optimal combination is worth it only if
  \[
    \Qave-\Qopt \;>\; \tfrac{m-1}{T}\,(\sigma_\varepsilon^2+\Qopt),
  \]
  which needs $T\gg m$ to hold (cost negligible); but in reality $T$ is of
  moderate size, so it often \textbf{fails}.

  \begin{alertblock}{The RSM departure}
    Use a \textbf{subset} of size $k\ll m$, with inflation
    $\dfrac{T-1}{T-k}\;\ll\;\dfrac{T-1}{T-m}$,\qquad
    i.e.\ cost $\dfrac{k-1}{T}\ll\dfrac{m-1}{T}$.
    \\[2pt]
    A window of \emph{net} gains opens even when full OLS is not worth estimating.
  \end{alertblock}
\end{frame}

%-----------------------------------------------------------------------------
\begin{frame}{What RSM adds over the literature}
  \begin{description}
    \item[Elliott--Liao ] When gains are largest, the optimum is a fixed
      subset with equal within-subset weights. But: \emph{which} forecasts?
      \emph{how large} the subset? --- left open.
  \end{description}
  \textbf{RSM closes both gaps and goes further:}
  \begin{enumerate}
    \item \textbf{Selection via randomisation} --- replaces an intractable
      $\binom{m}{k}$ search with uniform random draws + averaging.
      \item \textbf{In population}, RSM contains both equal-weighting ($k=1$) and optimal weights $k=m$.
    \item \textbf{A theory for $k^*$} --- the (potential) optimal subset size, which
      neither \citet*{ElliottLiao2025} nor \citet*{ChenMaung2023} provide.
  \end{enumerate}
  \begin{block}{The central question}
    In a panel of $m$ forecasters, \emph{how many} should we randomly
    subsample, and \emph{how} should we combine them, to minimise
    out-of-sample loss?
  \end{block}
\end{frame}

%=============================================================================
\section{Setup and Loss Function}
%=============================================================================

\begin{frame}{Setup}
  $m$ individual forecasts $f_t=(f_{1t},\dots,f_{mt})'$ of a scalar target
  $y_t=\mu_t+\epsilon_t$. Forecast errors: $u_t=f_t-\mu_t\im$.

  \begin{assumptionX}[Target innovation]
    $\E(\epsilon_t)=0,\quad \Var(\epsilon_t)=\sigma_\varepsilon^2>0,\quad
     \E(\epsilon_t u_{it})=0\ \ \forall i.$
  \end{assumptionX}
  \begin{assumptionX}[Forecast-error covariance]
    $\E(u_t)=0,\quad \Var(u_t)=\Sig\succ 0.$
  \end{assumptionX}

  Following \citet*{ElliottLiao2025}:
  \[
    \begin{pmatrix} y_t \\ f_t \end{pmatrix}
    \sim
    \left(
    \begin{pmatrix}\mu_t \\ \mu_t \end{pmatrix},
    \begin{pmatrix}\sigma_\varepsilon^2 & 0 \\ 0 & \Sig\end{pmatrix}
    \right).
  \]
  $\Sig$ summarises cross-forecast dispersion (diagonal: marginal error
  variances; off-diagonal: co-movement). $\sigma_\varepsilon^2$ is the
  \emph{irreducible noise} shared by every combination.
\end{frame}

%-----------------------------------------------------------------------------
\begin{frame}{The loss function}
  A combined forecast $w_m' f_t$ with $w_m'\im=1$ gives error
  $w_m'u_t-\epsilon_t$ and MSE
  \[
    L=\underbrace{(1-w_m'\im)^2\mu_t^2}_{\text{Bias}}
      +\underbrace{\sigma_\varepsilon^2+w_m'\Sig w_m}_{\text{Variance}}.
  \]
  Under the sum-to-one constraint the combination is \textbf{unbiased}, so
  the bias term vanishes:
  \[
    L=\sigma_\varepsilon^2+w_m'\Sig w_m .
  \]
  \begin{itemize}
    \item Comparison across methods reduces to minimising
      $w_m'\Sig w_m$.
    \item $\sigma_\varepsilon^2$ is common to all combinations --- but it
      \emph{enters} the loss and gets inflated by estimation error, so
      (unlike \citealp*{BootNibbering2019}) we keep it.
  \end{itemize}
\end{frame}

%-----------------------------------------------------------------------------
\begin{frame}{The two benchmarks}
  \begin{columns}[T]
    \begin{column}{0.5\textwidth}
      \begin{block}{Equal weighting}
        \[
          \wave=m^{-1}\im
        \]
        \[
          \Lave=\sigma_\varepsilon^2+\underbrace{m^{-2}\im'\Sig\im}_{\Qave}
        \]
        No weights estimated $\Rightarrow$ \emph{no} estimation penalty.
      \end{block}
    \end{column}
    \begin{column}{0.5\textwidth}
      \begin{block}{Optimal weighting}
        \[
          \wopt=\frac{\Sig^{-1}\im}{\im'\Sig^{-1}\im}
        \]
        \[
          \Lopt\approx(\sigma_\varepsilon^2+\Qopt)\,\frac{T-1}{T-m}
        \]
        with $\Qopt=(\im'\Sig^{-1}\im)^{-1}$.
      \end{block}
    \end{column}
  \end{columns}
  \vspace{4pt}
  \begin{alertblock}{Inflation factor}
    We use the \textbf{exact} out-of-sample inflation
    $I(T,k)=\dfrac{T-1}{T-k}$ for $k-1$ estimated coefficients.
    \citet*{ElliottLiao2025} use its first-order approximation
    $1+\tfrac{m-1}{T}$ (valid for $m\ll T$).
  \end{alertblock}
\end{frame}

%=============================================================================
\section{RSM: Definition and the Bagging Bridge}
%=============================================================================

\begin{frame}{RSM: subset weights}
  Draw a subset $S\subset\{1,\dots,m\}$, $|S|=k$, uniformly without
  replacement from $\binom{m}{k}$ subsets. Let $\Sigk$ be the $k\times k$
  principal submatrix.

  \begin{block}{Within-subset optimal weights}
    \[
      w_S=\frac{\Sigk^{-1}\ik}{\ik'\Sigk^{-1}\ik}.
    \]
  \end{block}
  Embed into $\mathbb{R}^m$ via the zero-padding map $P_S$:
  \[
    v_S=P_S w_S,\qquad \sum_{i=1}^m v_{S,i}=1 .
  \]
  Zero-padding preserves the quadratic form:
  \[
    v_S'\Sig v_S=w_S'\Sigk w_S=Q(S),\qquad
    Q(S)=(\ik'\Sigk^{-1}\ik)^{-1}.
  \]
\end{frame}

%-----------------------------------------------------------------------------
\begin{frame}{Three population risks}
  \begin{definitionX}[Average-subset RSM]
    Single random draw, expected loss over all subsets:
    \[
      \Qsub=\E_S[v_S'\Sig v_S]=\E_S[Q(S)].
    \]
  \end{definitionX}
  \begin{definitionX}[Finite-bagged RSM]
    $G$ independent draws, $\hat v_{k,G}=\tfrac1G\sum_g v_{S_g}$, then
    \[
      \QbagG=\E[\hat v_{k,G}'\,\Sig\,\hat v_{k,G}].
    \]
    \footnotesize Loss of the \emph{aggregated} forecast --- not the average of losses.
  \end{definitionX}
  \begin{definitionX}[Infinite-bagged RSM]
    $\bar v_k=\E_S[v_S]$ as $G\to\infty$, with deterministic risk
    \[
      \Qbag=\bar v_k'\Sig\bar v_k.
    \]
  \end{definitionX}
\end{frame}

%-----------------------------------------------------------------------------
\begin{frame}{Convexity: bagging never hurts}
  Since $\Sig\succ 0$, the map $v\mapsto v'\Sig v$ is strictly convex.
  Jensen's inequality gives
  \[
    \Qsub=\E_S[v_S'\Sig v_S]
    \;\ge\;
    (\E_S[v_S])'\Sig(\E_S[v_S])=\Qbag .
  \]

  \begin{itemize}
    \item Equality \emph{iff} $v_S$ is constant across subsets ---
      which happens only when $k=m$.
    \item Otherwise the inequality is \textbf{strict}: infinite bagging
      strictly reduces population risk relative to a single draw.
  \end{itemize}

  \begin{exampleblock}{Takeaway}
    Averaging weight vectors (bagging) is a free variance reduction ---
    no distributional assumptions needed beyond $\Sig\succ0$.
  \end{exampleblock}
\end{frame}

%-----------------------------------------------------------------------------
\begin{frame}{The finite-bagging bridge}
  \begin{propositionX}[Finite-bagging bridge]
    For independent subset draws,
    \[
      \QbagG=\Qbag+\frac{1}{G}\bigl(\Qsub-\Qbag\bigr),
    \]
    hence
    \[
      \Qbag\;\le\;\QbagG\;\le\;\Qsub .
    \]
  \end{propositionX}
  \begin{itemize}
    \item Holds for \emph{any} distribution of $\{f_t\}$.
    \item Increasing $G$ \textbf{monotonically} lowers the finite-bagged
      risk toward $\Qbag$.
    \item Special cases: $\QbagG=\Qsub$ at $G=1$; $\QbagG\to\Qbag$ as
      $G\to\infty$.
    \item Finite bagging is the practitioner's object: intermediate between
      the two limits, arbitrarily close to $\Qbag$ for large $G$.
  \end{itemize}
\end{frame}

%-----------------------------------------------------------------------------
\begin{frame}{What can we say with \emph{no} structure?}
  \footnotesize What survives if we know \emph{nothing} about $\Sig$
  beyond $\Sig\succ0$ --- no common loading, no distribution of $\{f_t\}$?
  \normalsize
  \begin{block}{Two-sided sandwich}
    \[
      \Qopt\;\le\;\Qbag\;\le\;\Qsub .
    \]
  \end{block}
  \begin{itemize}
    \item \textbf{Lower bound (feasibility).} $\bar v_k=\E_S[v_S]$ is an
      admissible sum-to-one combination ($\bar v_k'\im=1$). Since $\wopt$
      \emph{minimises} $w'\Sig w$ over all such $w$,
      $\;\Qbag=\bar v_k'\Sig\bar v_k\ge\Qopt$.
    \item \textbf{Upper bound (Jensen).} Convexity of $v\mapsto v'\Sig v$
      gives $\Qbag\le\Qsub$.
    \item \textbf{Anchored at the benchmarks.}
      $Q_1^{\mathrm{bag},\infty}=\Qave$, $\;Q_m^{\mathrm{bag},\infty}=\Qopt$.
  \end{itemize}
  \begin{alertblock}{Takeaway}
    In \emph{population} bagged RSM never beats the (infeasible) optimum and
    never does worse than a single draw. Any out-of-sample \emph{win over
    optimal} comes from the smaller inflation factor $\tfrac{T-1}{T-k}$ ---
    not from population risk.
  \end{alertblock}
\end{frame}

%=============================================================================
\section{RSM as a Unifier}
%=============================================================================

\begin{frame}{Boundary cases: $k=1$ and $k=m$}
  \textbf{Singletons ($k=1$):} every draw assigns full weight to one
  forecast, $v_{\{i\},i}=1$. Drawn with prob.\ $1/m$:
  \[
    \bar v_{1,i}=\E_S[v_{S,i}]=\tfrac1m\Rightarrow
    \bar v_1=\tfrac1m\im=\wave .
  \]
  \textbf{Full set ($k=m$):} only one subset, no randomness:
  \[
    \bar v_m=\frac{\Sig^{-1}\im}{\im'\Sig^{-1}\im}=\wopt .
  \]
  \begin{propositionX}[RSM as a unifier]
    Both benchmarks are special cases of RSM: $k=1$ collapses (in
    population) to \emph{equal weighting}; $k=m$ yields \emph{optimal
    weighting}. RSM interpolates as $k$ grows.
  \end{propositionX}
  \footnotesize Homogeneous forecasts $\Rightarrow$ small $k$ optimal;
  $T\gg m \Rightarrow k=m$. Most real panels are heterogeneous with short
  $T$ $\Rightarrow$ an interior $1<k<m$ wins.
\end{frame}

%-----------------------------------------------------------------------------
\begin{frame}{Global vs.\ conditional weights}
  With uniform size-$k$ subsets, $P(i\in S)=k/m$, so the global weight is
  \[
    \bar v_{k,i}=\E_S[v_{S,i}]
    =P(i\in S)\,\E_S[w_{S,i}\mid i\in S]
    =\frac{k}{m}\,\E_S[w_{S,i}\mid i\in S].
  \]
  \begin{propositionX}[Global--conditional correspondence]
    \[
      \bar v_{k,i}>\frac1m
      \quad\Longleftrightarrow\quad
      \E_S[w_{S,i}\mid i\in S]>\frac1k .
    \]
  \end{propositionX}
  \begin{itemize}
    \item A forecast earns above-average \emph{global} weight iff it earns
      above-average weight \emph{within the subsets it appears in}.
    \item ``Global importance is achieved locally.''
    \item Propositions on this slide and the last require \emph{only}
      $\Sig\succ0$ --- exact identities, no structural assumptions.
  \end{itemize}
\end{frame}

%=============================================================================
\section{Common-Loading Structure}
%=============================================================================

\begin{frame}{Common-loading heteroskedastic errors}
  \begin{assumptionX}[Forecast errors]
    \[
      u_t=\lambda F_t\im+\xi_t,\qquad
      \Var(\lambda F_t)=q>0,\quad
      \Var(\xi_t)=D=\mathrm{diag}(\sigma_1^2,\dots,\sigma_m^2).
    \]
  \end{assumptionX}
  A common shock (same loading $\lambda$) plus idiosyncratic noise:
  \[
    \Sig=q\,\im\im'+D,\qquad
    \Sigma_{e,S}=q\,\ik\ik'+D_S .
  \]
  \begin{definitionX}[Precision and precision sums]
    \[
      \tau_i=\sigma_i^{-2},\qquad
      A_S=\sum_{i\in S}\tau_i,\qquad
      A_m=\sum_{i=1}^m\tau_i .
    \]
  \end{definitionX}
\end{frame}

%-----------------------------------------------------------------------------
\begin{frame}{Subset risk under common loading}
  \begin{propositionX}[Subset risk]
    For any subset $S$ of size $k$,
    \[
      Q(S)=q+\frac{1}{A_S},
      \qquad
      w_{i,S}=\frac{\tau_i}{A_S}\ \ (i\in S).
    \]
  \end{propositionX}
  \begin{itemize}
    \item Loss $=$ irreducible common component $q$ $+$ inverse sum of
      drawn precisions.
    \item $Q(S)$ is strictly decreasing in $k$ along nested subsets ($\tau_i>0$);
      minimised at $k=m$ ($\to\Qopt$).
    \item \textbf{But} the inflation factor $I(T,m)$ may stop full optimal
      from winning in practice.
    \item Optimal subset weights depend only on the precisions $\tau$ ---
      \emph{not} on $q$.
  \end{itemize}
\end{frame}

%-----------------------------------------------------------------------------
\begin{frame}{Optimal vs.\ equal weighting: the $B$--$H$ decomposition}
  Optimal weight $w_i^{\mathrm{opt}}=\tau_i/A_m$, loss $\Qopt=q+1/A_m$.
  Equal weighting ignores heterogeneity:
  $\Qave=q+\bigl(\sum_i\tau_i^{-1}\bigr)/m^2$.

  Define $B=1/\bar\tau$ (reciprocal mean precision; idiosyncratic
  \emph{scale}) and $H=\overline{\sigma^2}\,\bar\tau$ (arithmetic/harmonic
  variance ratio; \emph{dispersion}, $H\ge1$). Then:
  \[
    \boxed{\;\Qopt=q+\frac{B}{m},\qquad
    \Qave=q+\frac{BH}{m},\qquad
    \Qave-\Qopt=\frac{B(H-1)}{m}\;}
  \]
  \begin{itemize}
    \item Equal weighting can never beat optimal in population.
    \item They coincide iff variances are homogeneous ($H=1$).
    \item RSM with $k>1$ pays off only when variances are
      \textbf{dispersed enough} ($H$ large).
  \end{itemize}
\end{frame}

%=============================================================================
\section{RSM Risk and the Optimal $k^*$}
%=============================================================================

\begin{frame}{Average-subset RSM and the square-root rule}
  \begin{propositionX}[Average-subset risk]
    \[
      \Qsub=q+\frac{1}{k\bar\tau}
        +\frac{m-k}{k^2(m-1)}\frac{v_\tau}{\bar\tau^3}+O(k^{-3})
      \;=\;q+\frac{B}{k}+O(k^{-2}).
    \]
  \end{propositionX}
  Same shape as $\Qopt$, but divided by $k$ instead of $m$ --- and the
  inflation cost is smaller, $I(T,k)<I(T,m)$.
  \begin{block}{Optimal $k$ for average-subset RSM}
    \[
      \Lsub\approx\Bigl(\sigma_\varepsilon^2+q+\tfrac{B}{k}\Bigr)\frac{T-1}{T-k}\to\min
      \;\Longrightarrow\;
      k_{\mathrm{sub}}^*=\frac{-B+\sqrt{B^2+(\sigma_\varepsilon^2+q)BT}}{\sigma_\varepsilon^2+q}.
    \]
    \[
      k^*=\min\{m,\max\{1,\mathrm{round}(k_{\mathrm{sub}}^*)\}\}
      \;\sim\;O(\sqrt{T}).
    \]
  \end{block}
\end{frame}

%-----------------------------------------------------------------------------
\begin{frame}{Why the bagged optimum resists the same derivation}
  The $\sqrt{T}$ rule came from Taylor-expanding an \emph{unconditional}
  expectation ($A_S=k\bar\tau+\epsilon_S$, $\E_S[\epsilon_S]=0$, known variance):
  \[
    \Qsub=q+\E_S\!\Bigl[\tfrac{1}{A_S}\Bigr]
      =q+\frac{1}{k\bar\tau}+\frac{\Var_S(A_S)}{(k\bar\tau)^3}+\cdots
  \]
  The bagged risk instead needs the \emph{conditional} object
  \[
    g_i(k)=k\,\E_S\!\Bigl[\tfrac{1}{A_S}\,\Big|\,i\in S\Bigr]
      =k\,\E\!\Bigl[\frac{1}{\tau_i+A_{S\setminus\{i\}}}\Bigr].
  \]
  \footnotesize
  \begin{itemize}
    \item A \textbf{shifted conditional harmonic mean}: $A_S=\tau_i+$ (sum of
      $k-1$ \emph{other} precisions).
    \item Depends on the size-$(k-1)$ precision sums of $\{\tau_j\}_{j\ne i}$
      --- \textbf{a different distribution for each $i$}.
    \item No closed form without restricting $\{\tau_i\}$: conditioning on
      $i\in S$ destroys the factorisation.
  \end{itemize}
  \normalsize
  \begin{alertblock}{Way forward}
    Characterise $\Rk$ by Monte Carlo --- yet its \emph{endpoints} are exact,
    enough to compare bagged RSM to both benchmarks with \emph{no} functional form.
  \end{alertblock}
\end{frame}

%-----------------------------------------------------------------------------
\begin{frame}{Infinite-bagged weights and the conservation law}
  Let $g_i(k)=k\,\E_S\!\bigl[\tfrac{1}{A_S}\mid i\in S\bigr]$. The
  infinite-bagged weight on forecast $i$ is
  \[
    \bar v_{k,i}=\E_S[v_{S,i}]
      =\frac{k}{m}\,\E_S\!\Bigl[\frac{\tau_i}{A_S}\,\Big|\,i\in S\Bigr]
      =\frac{\tau_i\,g_i(k)}{m}.
  \]
  Since $\sum_{i\in S}w_{S,i}=1$ for every $S$, we have
  $\sum_i\bar v_{k,i}=1\Rightarrow\sum_i\tau_i g_i(k)=m$ for all $k$.
  \begin{propositionX}[Conservation law]
    For every $k\in\{1,\dots,m\}$,
    \[
      \sum_{i=1}^m \tau_i\Bigl[g_i(k)-\frac{1}{\tau_i}\Bigr]=0 .
    \]
  \end{propositionX}
  \footnotesize
  \begin{itemize}
    \item $1/\tau_i$ is the $g$-reference giving \emph{equal} weights:
      at $k=1$, $g_i(1)=1/\tau_i$ and $\bar v_{1,i}=1/m$.
    \item Deviation $\tau_i[g_i(k)-1/\tau_i]=m\bar v_{k,i}-1$ is the
      \emph{excess weight above equal weighting}, precision-rescaled.
    \item RSM redistributes weight relative to EW with \textbf{zero}
      net precision-weighted gain --- holds for all $k$, all $\{\tau_i\}$,
      no distributional assumption.
  \end{itemize}
\end{frame}

%-----------------------------------------------------------------------------
\begin{frame}{Exact excess risk of the infinite-bagged RSM}
  Weight deviation $\Delta_{i,k}=\bar v_{k,i}-w_i^{\mathrm{opt}}$ and excess
  risk $\Rk=\Qbag-\Qopt$.
  \begin{propositionX}[Exact excess risk]
    Under common loading,
    \[
      \Rk=\sum_{i=1}^m\frac{\bigl(\bar v_{k,i}-w_i^{\mathrm{opt}}\bigr)^2}{\tau_i}
        \;=\;\frac{1}{m^2}\sum_{i=1}^m\tau_i\Bigl(g_i(k)-\frac{1}{\bar\tau}\Bigr)^{\!2}\;\ge0 .
    \]
  \end{propositionX}
  \footnotesize
  The common-loading term drops because $\im'\Delta_k=0$; only the
  idiosyncratic part $\sum_i\sigma_i^2\Delta_{i,k}^2$ survives.
  \begin{itemize}
    \item \textbf{Endpoints (exact):} $\Rk=0$ at $k=m$ ($\bar v_{m}=\wopt$);
      $R_1=\Qave-\Qopt=\dfrac{B(H-1)}{m}>0$ unless all $\tau_i$ equal.
    \item \textbf{Sharp Jensen bound:}
      $\Rk\le\Qsub-\Qopt=B\,z_k$, with $z_k=\tfrac1k-\tfrac1m$ ---
      tight iff homogeneous.
    \item $\Rk$ \textbf{strictly decreasing} in $k$ --- its decay pins down $k^*$.
  \end{itemize}
\end{frame}

%-----------------------------------------------------------------------------
\begin{frame}{Comparing bagged RSM to benchmarks --- with \emph{no} form}
  Write the inflated bagged loss using only its \emph{endpoints} and sign:
  \[
    \Lbag=\bigl(A+\Rk\bigr)\,\mathcal{I}(T,k),\qquad
    A:=\sigma_\varepsilon^2+\Qopt,\quad \mathcal{I}(T,k)=\tfrac{T-1}{T-k}.
  \]
  \footnotesize Structural facts only: $R_1=\tfrac{B(H-1)}{m}$, $R_m=0$,
  $\Rk\ge0$, and $\Rk\downarrow$ in $k$.\normalsize
  \begin{propositionX}[Bagged RSM vs.\ benchmarks, $1\le k<m<T$]
    \begin{enumerate}
      \item[(i)] \textbf{Population bracketing:}
        $\Qopt\le\Qbag\le\Qave$ for every $k$.
      \item[(ii)] \textbf{Endpoint domination:}
        $L_1^{\mathrm{bag},\infty}=\Lave$ and
        $L_m^{\mathrm{bag},\infty}=\Lopt$, hence
        \[
          \min_{1\le k\le m}\Lbag\;\le\;\min\bigl(\Lave,\,\Lopt\bigr).
        \]
    \end{enumerate}
  \end{propositionX}
  \footnotesize
  (ii) is \textbf{unconditional} --- it needs neither monotonicity nor a
  closed form, only that $k=1$ and $k=m$ are admissible. RSM is therefore
  \emph{never worse} than the better of the two benchmarks.
\end{frame}

%-----------------------------------------------------------------------------
\begin{frame}{Exact comparison conditions (still no form)}
  \begin{propositionX}[Continued]
    \begin{enumerate}
      \item[(iii)] \textbf{Exact comparison conditions:}
        \[
          \Lbag<\Lopt\iff \Rk<\frac{A\,(m-k)}{T-m},
        \]
        \[
          \Lbag<\Lave\iff A\,(k-1)<R_1(T-k)-\Rk(T-1).
        \]
      \item[(iv)] \textbf{Interior strict domination:} the interior optimum
        beats \emph{both} endpoints provided
        \[
          R_1(T-2)-R_2(T-1)>A
          \quad\text{and}\quad
          R_{m-1}<\frac{A}{T-m}.
        \]
    \end{enumerate}
  \end{propositionX}
  \begin{block}{Conservative (tractable) version}
    \footnotesize Jensen gives $\Rk\le B\,z_k$; substituting $B z_k$ for $\Rk$
    in (iii) yields \emph{sufficient} conditions --- exactly those from the
    average-subset $\sqrt{T}$ analysis.
  \end{block}
\end{frame}

%-----------------------------------------------------------------------------
\begin{frame}{Choosing $k^*$: when is it small, large, or interior?}
  \renewcommand{\arraystretch}{1.35}
  \begin{center}
  \begin{tabular}{@{}p{0.36\textwidth} p{0.30\textwidth} c@{}}
    \toprule
    \textbf{Regime} & \textbf{Condition} & \textbf{$k^*$} \\
    \midrule
    Homogeneous forecasts & $H\approx1$, so $R_1\approx0$ &
      $\approx 1$ \\
    Long sample / small pool & $T\gg m$ (inflation negligible) &
      $\approx m$ \\
    Heterogeneous, moderate $T$ & $H$ large, $T\sim m$ &
      interior $1<k^*<m$ \\
    Endogenous / $\Sig$ unknown & estimated forecasts &
      $\approx\lfloor\sqrt{T}\rfloor$ \\
    \bottomrule
  \end{tabular}
  \end{center}
  \begin{itemize}
    \item \textbf{Homogeneous} $\Rightarrow$ gains vanish, RSM collapses to
      \emph{equal weighting}.
    \item \textbf{$T\gg m$} $\Rightarrow$ estimation cheap, RSM tends to
      \emph{full optimal}.
    \item \textbf{Dispersed errors, short $T$} $\Rightarrow$ an interior
      $k^*$ strictly dominates both benchmarks (Prop., part iv).
    \item \textbf{Endogenous case} (forecasts themselves estimated):
      $k^*\approx\lfloor\sqrt{T}\rfloor$ is a robust \emph{starting point}.
  \end{itemize}
\end{frame}

%-----------------------------------------------------------------------------
\begin{frame}{Bounding $R_k$ from both sides: the optimal-$k$ window}
  The Jensen ceiling $\Rk\le Bz_k$ has a matching \emph{projection floor},
  giving an exact two-sided bracket:
  \[
    \underbrace{\frac{V_\delta^2}{m(H-1)\bar\tau}\,z_k^2}_{\text{projection floor}}
    \;\lesssim\;\Rk\;\lesssim\;
    \underbrace{B\,z_k}_{\text{Jensen ceiling}},
    \qquad V_\delta=\frac{v_\tau}{\bar\tau^2},\ \ z_k=\tfrac1k-\tfrac1m .
  \]
  {\footnotesize (Floor: Cauchy--Schwarz projection of $\bar v_k-\wopt$ onto $\wave-\wopt$; exact for $\le2$ distinct precisions.)}
  \begin{block}{A bounded window for $k^*$}
    $0<z_k<1\Rightarrow z_k^2<z_k$, so $\Rk\sim z_k^\alpha$ with
    $\alpha\in[1,2]$, and the optimal subset size is \textbf{confined} to
    $\;\Theta(T^{1/3})\le k^*\le\Theta(T^{1/2}).$
  \end{block}
  {\footnotesize
  \textbf{Cube root} ($\alpha=2$): bounded coherence \emph{and} dispersed
  precisions.\quad
  \textbf{Square root} ($\alpha=1$): a dominant forecaster, or the endogenous
  case (estimated weights).}
  \vspace{2pt}

  \centering\structure{No functional form needed --- Monte Carlo locates the optimum inside the window.}
\end{frame}

%=============================================================================
\begin{frame}{Summary of Section 4}
  \begin{enumerate}
    \item Unbiased combinations $\Rightarrow$ minimise variance
      $\sigma_\varepsilon^2+w'\Sig w$.
    \item RSM \textbf{unifies} equal weighting ($k=1$) and optimal weighting
      ($k=m$); exact identities, no structure needed.
    \item \textbf{Bagging bridge:}
      $\Qbag\le\QbagG\le\Qsub$, monotone in $G$; with \emph{no} structure,
      $\Qopt\le\Qbag\le\Qsub$.
    \item Under common loading: $Q(S)=q+1/A_S$, the
      $\Qave-\Qopt=B(H-1)/m$ decomposition, and the \textbf{conservation law}
      $\sum_i\tau_i[g_i(k)-1/\tau_i]=0$.
    \item \textbf{Exact excess risk}
      $\Rk=\sum_i(\bar v_{k,i}-w_i^{\mathrm{opt}})^2/\tau_i\ge0$;
      square-root rule for average-subset $k^*$ and a general FOC balancing
      gain vs.\ inflation cost.
  \end{enumerate}
  \vfill
  \centering\structure{\large Next: Monte Carlo evidence (Section 5).}
\end{frame}

%=============================================================================
\begin{frame}[allowframebreaks]{References}
  \footnotesize
  \begin{thebibliography}{Elliott et al.(2013)}
    \bibitem[Bates and Granger(1969)]{BatesGranger1969} Bates, J.M.\ \& Granger, C.W.J.\ (1969).
      The combination of forecasts. \emph{OR}.
    \bibitem[Granger and Ramanathan(1984)]{GrangerRamanathan1984} Granger, C.W.J.\ \& Ramanathan, R.\ (1984).
      Improved methods of combining forecasts. \emph{J.\ Forecasting}.
    \bibitem[Elliott and Liao(2025)]{ElliottLiao2025} Elliott, G.\ \& Liao, Y.\ (2025).
      Forecast combination and the combination puzzle.
    \bibitem[Elliott et al.(2013)]{ElliottGarganoTimmermann2013} Elliott, G., Gargano, A.\
      \& Timmermann, A.\ (2013). Complete subset regressions.
      \emph{J.\ Econometrics}.
    \bibitem[Chen and Maung(2023)]{ChenMaung2023} Chen, J.\ \& Maung, K.\ (2023).
      Nonparametric time-varying forecast combination.
    \bibitem[Boot and Nibbering(2019)]{BootNibbering2019} Boot, T.\ \& Nibbering, D.\ (2019).
      Forecasting using random subspace methods. \emph{J.\ Econometrics}.
  \end{thebibliography}
\end{frame}

\end{document}
